% !TEX root = ./thesis.tex
% =========================================================================
% CHAPTER 2: LITERATURE REVIEW
% =========================================================================
\chapter{Literature Review}

\section{Introduction and Evolution of Cuttings Transport Research}
The study of drill cuttings transport in deviated and horizontal wellbores has evolved across three distinct scientific eras over the past four decades \citep{Li2026, Khalili2025}. As synthesized in the comprehensive bibliometric review by \citet{Li2026}, the first era (1980s--1990s) focused on foundational laboratory flow-loop experiments aimed at identifying critical transport velocities, inclination effects, and empirical correlations \citep{Tomren1986, Sanchez1999}. While these early empirical correlations provided rapid field screening tools, they were inherently restricted to the specific pipe diameters, fluid types, and inclination ranges of the test facilities from which they were regressed, resulting in large prediction errors when extrapolated to different wellbore geometries \citep{Sorgun2010, Li2026}.

The second era (1990s--2000s) introduced one-dimensional mechanistic two-layer and three-layer models \citep{Gavignet1989, Kamp1999, Hyun2000}. These models divided the horizontal wellbore cross-section into a stationary or moving porous cuttings bed at the bottom, an intermediate dispersing layer, and a clear or dilute suspended fluid layer at the top, applying mass and momentum conservation across each layer \citep{Sorgun2010, Li2026}. However, as critically noted by \citet{Kamp1999} and \citet{Sorgun2010}, 1D layered models rely on simplified cross-sectional averaging and struggle to capture the complex, three-dimensional secondary swirl, asymmetric shear stress distribution, and localized bed erosion induced by drillpipe rotation in an eccentric annulus.

The third and current era (2010--present) is defined by high-fidelity three-dimensional Computational Fluid Dynamics (CFD) and coupled CFD--Discrete Element Method (CFD--DEM) simulations \citep{Shehab2015, Hajipour2020, Khaled2020, Li2026}. By solving the conservation equations on finely discretized 3D meshes, CFD resolves local velocity deficits in the narrow eccentric clearance, interfacial bed shear stresses, and phase volume fraction distributions ($\alpha_s$) that cannot be observed experimentally \citep{Li2026}. In this chapter, the foundational experimental and numerical literature is critically reviewed to establish a rigorous justification for how every geometric, physical, rheological, and numerical parameter in the present study is configured, while identifying the specific research gaps that motivate our open-source OpenFOAM methodology.

\section{Physics of Solid--Liquid Flow and Bed Regimes in Horizontal Annuli}
In horizontal and near-horizontal wellbores ($60^\circ$ to $90^\circ$ inclination from vertical), the radial component of gravity drives solid rock particles toward the lower annular wall immediately upon entering the flow stream \citep{Tomren1986, Osgouei2010}. Whether a settled particle remains trapped on the bottom wall or is re-entrained into the axial flow depends on the local force balance between stabilizing gravitational/buoyancy forces and mobilizing hydrodynamic drag, shear lift, and mechanical agitation forces \citep{Khalili2025, Li2026}.

Based on extensive experimental high-speed camera observations in horizontal flow loops \citep{Tomren1986, Osgouei2010}, the solid--liquid mixture in a horizontal eccentric annulus transitions through three primary macroscopic flow regimes as hydraulic and mechanical energy increase:
\begin{enumerate}
    \item \textbf{Stationary Cuttings Bed Regime:} At low superficial fluid velocities (typically $v_l < 1.2\text{ m/s}$) and zero or low pipe rotation, the hydrodynamic drag and near-bed lift forces fall below the critical threshold required for incipient particle motion \citep{Osgouei2010, Li2026}. Cuttings settle into a tightly packed, immobile granular bed along the lower casing floor, reaching solid volume fractions approaching the maximum packing limit ($\alpha_s \approx 0.55\text{--}0.63$). Fluid bypasses the bed entirely through the upper open channel, resulting in severe flow restriction and high stuck-pipe risk \citep{Shehab2015, Li2026}.
    \item \textbf{Moving and Saltating Bed Regime:} As the superficial fluid velocity increases into the intermediate transition zone ($v_l \approx 1.2\text{ to } 1.6\text{ m/s}$) or moderate drillpipe rotation is applied, interfacial shear stress exceeds the critical bed erosion threshold \citep{Osgouei2010, Khalili2025}. Surface particles roll, slide, or saltate (bounce) along the top of the bed in the form of migrating dunes or a continuously sliding granular layer, progressively reducing the stationary bed height \citep{Osgouei2010, Li2026}.
    \item \textbf{Asymmetric Suspension and Dispersed Flow Regime:} At high axial velocities ($v_l > 1.8\text{--}2.0\text{ m/s}$) coupled with vigorous drillpipe rotation ($\Omega \ge 80\text{--}120\text{ RPM}$), turbulent dispersion, hydrodynamic lift, and rotational tangential sweeping overcome gravitational settling \citep{Osgouei2010, Shehab2015, Khalili2025}. The compacted bed is completely eroded, and solid particles are transported as an asymmetric or heterogeneous dispersed suspension across the annular cross-section \citep{Osgouei2010, Li2026}.
\end{enumerate}

Capturing the exact operational boundaries between these three regimes is essential for hole-cleaning design and directly motivates Objective 3 of the present study.

\section{Critical Literature Review and Justification of Simulation Parameters}
A numerical simulation is only as reliable as the physical realism of its input parameters and its linkage to verified experimental benchmarks \citep{Li2026}. Rather than selecting arbitrary geometry or flow conditions, every parameter in the present OpenFOAM framework is systematically anchored to established experimental flow-loop studies and peer-reviewed CFD benchmarks.

\subsection{Justification of Annular Geometry and Eccentricity ($D_w, D_p, L, e = 0.623$)}
In real horizontal drilling operations, the drillstring never remains concentric within the borehole; gravity pulls the heavy steel pipe toward the bottom wall \citep{Sorgun2010, Saad2024}. Dimensionless pipe eccentricity ($e$) is defined as:
\begin{equation}
    e = \frac{2\delta}{D_w - D_p}
\end{equation}
where $\delta$ is the radial offset distance between the centerline of the outer wellbore ($D_w$) and the centerline of the inner drillpipe ($D_p$) \citep{Sorgun2010, Shehab2015}. An eccentricity of $e = 0$ corresponds to a perfectly concentric annulus, whereas $e = 1.0$ represents a fully eccentric pipe in direct line contact with the bottom wall.

To establish a rigorous validation baseline, this study adopts the exact geometric configuration of the Middle East Technical University (METU) horizontal flow-loop facility documented in the Ph.D. dissertations of \citet{Osgouei2010} and \citet{Sorgun2010}, and subsequently modeled in ANSYS CFX by \citet{Shehab2015} and \citet{Ofei2015}. Specifically:
\begin{itemize}
    \item \textbf{Outer Wellbore Diameter ($D_w$):} Set to $0.0739\text{ m}$ ($2.91\text{ in}$), representing a scaled slim-hole horizontal test section \citep{Osgouei2010, Shehab2015}.
    \item \textbf{Inner Drillpipe Diameter ($D_p$):} Set to $0.0470\text{ m}$ ($1.85\text{ in}$), yielding a diameter ratio of $\kappa = D_p/D_w = 0.636$ and a hydraulic diameter of $D_h = D_w - D_p = 0.0269\text{ m}$ \citep{Osgouei2010, Shehab2015}.
    \item \textbf{Eccentricity ($e = 0.623$):} The inner pipe is offset downward along the vertical axis by $\delta = 8.38\text{ mm}$, yielding a fixed positive eccentricity of $e = 0.623$ \citep{Osgouei2010, Shehab2015}. We specifically select $e = 0.623$ for two critical reasons: (1) it replicates the exact mechanical stand-off created by drillpipe tool joints resting on the lower wellbore wall in the \citet{Osgouei2010} experiments, and (2) unlike $e = 1.0$ which creates a zero-thickness singular contact line that causes severe mesh skewness and finite-volume divergence, $e = 0.623$ preserves a realistic narrow lower gap of $\sim 5.07\text{ mm}$ ($0.00507\text{ m}$) while expanding the upper gap to $\sim 21.83\text{ mm}$. This allows structured hexahedral boundary-layer cells to resolve the low-side creeping flow zone cleanly.
    \item \textbf{Axial Domain Length ($L = 3.3\text{ m}$):} As noted by \citet{Hajipour2020}, simulating an entire horizontal wellbore spanning thousands of meters is computationally intractable in 3D CFD. \citet{Osgouei2010} and \citet{Shehab2015} demonstrated that an axial test section of $L = 3.3\text{ m}$ ($L/D_h \approx 122.7$) is more than sufficient to eliminate inlet entrance effects and achieve hydrodynamically fully developed multiphase flow and steady cuttings bed thickness.
\end{itemize}

\subsection{Justification of Cuttings Phase Properties ($d_p, \rho_s, \text{ROP}$)}
The physical characteristics of the drilled solid phase directly govern the particle settling velocity, Stokes number, and interphase drag transfer \citep{Hajipour2020, Khaled2020}. 
\begin{itemize}
    \item \textbf{Particle Diameter ($d_p = 2.01\text{ mm}$) and Density ($\rho_s = 2761.4\text{ kg/m}^3$):} We adopt a uniform spherical particle diameter of $d_p = 0.00201\text{ m}$ and a solid material density of $\rho_s = 2761.4\text{ kg/m}^3$ directly from the sieve-analyzed sandstone cuttings used in the \citet{Osgouei2010} flow-loop experiments and the \citet{Shehab2015} CFD benchmark. Critically, \citet{Sorgun2010} and \citet{Hajipour2020} observed that medium-to-coarse cuttings around $2.0\text{ mm}$ are particularly prone to forming compact stationary beds in horizontal wells, making this particle size an ideal critical test case for evaluating bed erosion efficiency.
    \item \textbf{Rate of Penetration ($\text{ROP} = 60\text{ ft/hr}$):} The mass flux of solids entering the wellbore inlet is governed by the drilling Rate of Penetration, fixed at $\text{ROP} = 60\text{ ft/hr}$ ($0.00508\text{ m/s}$), matching the exact feed rate of \citet{Osgouei2010} and \citet{Shehab2015}. In the Eulerian framework, this corresponds to an inlet solid volume fraction of $\alpha_{s,\text{in}} \approx 0.04$ to $0.08$ depending on the carrier fluid flow rate, allowing one-to-one validation of steady-state annular cuttings concentration ($C_c\%$) and pressure drop ($\Delta P/L$).
\end{itemize}

\subsection{Justification of Operational Parameters: Fluid Velocity ($v_l$) and Pipe Rotation ($\Omega$)}
To systematically capture the transition from a hazardous stationary cuttings bed to a clean dispersed flow, the operational matrix must span across the critical hydrodynamic thresholds identified in literature:
\begin{itemize}
    \item \textbf{Superficial Fluid Inlet Velocity ($v_l = 1.2, 1.6, 2.0, 2.4\text{ m/s}$):} In their recent review of multiphase wellbore hydraulics, \citet{Khalili2025} highlighted that when annular fluid velocity falls below $1.5\text{ m/s}$, gravitational settling dominates and stationary cuttings beds rapidly form; once velocity surpasses $1.5\text{ m/s}$, axial drag and shear lift begin to dominate. Similarly, \citet{Osgouei2010} and \citet{Shehab2015} tested superficial velocities spanning $3.9\text{ to } 7.9\text{ ft/s}$ ($1.19\text{ to } 2.41\text{ m/s}$). Selecting $v_l = 1.2, 1.6, 2.0,\text{ and } 2.4\text{ m/s}$ strategically brackets this physical transition: $1.2\text{ m/s}$ captures the sub-critical stationary bed regime, $1.6\text{ m/s}$ captures the incipient moving-bed transition right above the $1.5\text{ m/s}$ threshold, and $2.0\text{--}2.4\text{ m/s}$ captures high-velocity turbulent suspension.
    \item \textbf{Drillpipe Rotational Speed ($\Omega = 0, 40, 80, 120\text{ RPM}$):} Early experimental work by \citet{Sanchez1999} and subsequent flow-loop investigations by \citet{Sorgun2010} and \citet{Osgouei2010} proved that mechanical rotation of the inner drillpipe significantly enhances cuttings removal in eccentric horizontal wells by imparting tangential Couette shear and secondary Taylor-like vortices that sweep cuttings from the narrow bottom gap into the high-velocity upper stream. However, \citet{Shehab2015} only validated two rotational speeds ($0\text{ and } 60\text{ RPM}$) in Newtonian water, while \citet{Khaled2021} noted that $0\text{ to } 200\text{ RPM}$ represents the critical operational range in field drilling, with diminishing returns and severe drillstring vibration risks occurring above $120\text{ RPM}$ \citep{Hajipour2020}. Therefore, evaluating $\Omega = 0, 40, 80,\text{ and } 120\text{ RPM}$ provides uniform $40\text{ RPM}$ increments that map the non-linear onset of rotational bed erosion across typical field rotary speeds \citep{Osgouei2010, Hajipour2020}.
\end{itemize}

\subsection{Justification of Non-Newtonian Drilling Mud Rheology ($\tau_y, K, n$)}
A major limitation in several foundational CFD studies—most notably \citet{Shehab2015}—is the simplification of the carrier liquid as pure Newtonian water ($\rho_l = 998.5\text{ kg/m}^3, \mu = 0.001\text{ Pa}\cdot\text{s}$). While using water is valuable for initial baseline validation against water-sand flow-loop runs in \citet{Osgouei2010}, real field drilling fluids are complex bentonite or polymer suspensions that exhibit strong non-Newtonian behavior \citep{Sorgun2010, Hajipour2020}.

As established by \citet{Kelessidis2006}, \citet{Hajipour2020}, and \citet{Khaled2020}, the three-parameter \textbf{Herschel--Bulkley rheological model} ($\tau = \tau_y + K\dot{\gamma}^n$) is the most accurate constitutive equation for water-based drilling muds because it simultaneously accounts for both a finite yield stress ($\tau_y$) at zero shear rate and shear-thinning power-law behavior ($n < 1$) at elevated shear rates. In contrast, the two-parameter Bingham Plastic model overpredicts viscosity at high shear rates, while the Ostwald--de Waele Power-Law modelused by \citet{Ferroudji2022} fails to capture the static yield stress ($\tau_y$) responsible for forming an unyielded gelled plug core at low shear rates \citep{Kelessidis2006, Hajipour2020}.

Based on the experimental viscometer measurements of water-based polymer muds in \citet{Sorgun2010}, \citet{Kelessidis2006}, and \citet{Hajipour2020}, our non-Newtonian fluid parameters are set to:
\begin{itemize}
    \item \textbf{Yield Stress ($\tau_y = 4.0\text{ Pa}$):} Equivalent to $\approx 8.35\text{ lbf}/100\text{ ft}^2$, falling squarely within the $1.0\text{ to } 8.5\text{ Pa}$ ($2\text{ to } 17\text{ lbf}/100\text{ ft}^2$) range reported by \citet{Hajipour2020} for low-to-medium yield-point drilling muds. Dividing by fluid density ($\rho_l = 1000\text{ kg/m}^3$) yields the OpenFOAM kinematic yield stress parameter $\tau_0 = \tau_y/\rho_l = 4.0 \times 10^{-3}\text{ m}^2/\text{s}^2$.
    \item \textbf{Consistency Index ($K = 0.4\text{ Pa}\cdot\text{s}^n$) and Flow Behavior Index ($n = 0.65$):} Matches the experimental polymer mud formulations tested by \citet{Sorgun2010} ($K = 0.289\text{--}0.806\text{ Pa}\cdot\text{s}^n$) and \citet{Hajipour2020} ($n = 0.61\text{--}0.68$), giving a kinematic consistency index of $k = K/\rho_l = 4.0 \times 10^{-4}\text{ m}^2/\text{s}^{2-n}$.
    \item \textbf{Regularization Viscosity Bounds ($\nu_0 = 1.0 \times 10^{-2}\text{ m}^2/\text{s}$):} In finite-volume solvers, as local shear rate approaches zero ($\dot{\gamma} \to 0$) in the wide upper annulus, the analytical apparent viscosity $\nu_{\text{eff}} = \tau_0/\dot{\gamma} + k\dot{\gamma}^{n-1}$ approaches infinity. Following OpenFOAM's regularized Herschel--Bulkley implementation, capping the maximum unyielded viscosity at $\nu_0 = 1.0 \times 10^{-2}\text{ m}^2/\text{s}$ ($\mu_0 = 10\text{ Pa}\cdot\text{s}$) accurately reproduces a rigid unyielded plug core while maintaining matrix diagonal dominance and numerical stability.
\end{itemize}

\subsection{Critique of Multiphase Modeling Approaches (Eulerian--Eulerian vs.\ CFD--DEM)}
Two primary numerical frameworks exist for simulating solid--liquid cuttings transport in 3D CFD: the Eulerian--Lagrangian Discrete Element Method (CFD--DEM) and the Eulerian--Eulerian Two-Fluid Model (TFM) \citep{Hajipour2020, Li2026}.

In CFD--DEM (e.g., \citealt{Akhshik2015}; reviewed in \citealt{Li2026}), the liquid is treated as a continuum while every individual rock cutting is tracked via Newton's second law of motion, resolving discrete particle--particle collisions directly. While CFD--DEM provides high particle-scale fidelity, \citet{Li2026} and \citet{Hajipour2020} emphasize that its computational cost scales prohibitively with particle count ($O(N_p \log N_p)$). In a $3.3\text{ m}$ horizontal wellbore containing a dense cuttings bed ($\alpha_s > 0.50$), tracking millions of $2\text{ mm}$ particles across a 16-run parametric matrix requires supercomputing clusters and weeks of runtime \citep{Li2026}.

Conversely, the \textbf{Eulerian--Eulerian Two-Fluid Model (TFM)} treats both the liquid mud and the solid cuttings as interpenetrating continua governed by phase-averaged Navier--Stokes equations \citep{Shehab2015, Hajipour2020, Khaled2020}. Particle--particle collisions, frictional packing stresses, and granular pressure ($P_s$) within the dense cuttings bed are closed mathematically using the \textbf{Kinetic Theory of Granular Flow (KTGF)} derived from the Boltzmann equation for granular gases \citep{Gidaspow1994, Syamlal1993, Hajipour2020}. As demonstrated by \citet{Hajipour2020} and \citet{Khaled2020}, the Eulerian--Eulerian KTGF framework coupled with the \citet{Wen1966} interphase drag closure achieves excellent macro-scale accuracy in predicting bed height, solid volume fraction ($\alpha_s$), and pressure loss ($\Delta P/L$) at a fraction of the computational cost of CFD--DEM, making it the optimal engineering choice for executing our 16-run parametric matrix in OpenFOAM (\texttt{twoPhaseEulerFoam}).

\section{Critical Evaluation of Benchmark Studies and Their Limitations}
To clearly position the present work within the body of knowledge, the primary benchmark studies used to guide our methodology are critically evaluated below in terms of what they achieved, how they motivate our approach, and why their work must be extended:

\begin{enumerate}
    \item \textbf{\citet{Osgouei2010} (Experimental Benchmark):} Conducted comprehensive horizontal flow-loop experiments ($D_w = 0.0739\text{ m}, D_p = 0.0470\text{ m}, e = 0.623, L = 3.3\text{ m}$) measuring steady-state frictional pressure drop ($\Delta P/L$) and average cuttings concentration ($C_c\%$). \textit{Why cited:} Provides the primary ground-truth experimental dataset used to validate our OpenFOAM model's baseline predictions. \textit{Limitation:} As an experimental study, it could only measure bulk volume-averaged quantities and could not resolve internal 3D velocity vectors, local effective viscosity gradients ($\nu_{\text{eff}}$), or cross-sectional shear stresses inside the narrow eccentric gap.
    
    \item \textbf{\citet{Shehab2015} \& \citet{Ofei2015} (ANSYS CFX Numerical Benchmark):} Replicated the \citet{Osgouei2010} geometry in commercial software (ANSYS 15.0 CFX) using an unstructured tetrahedral/prism mesh ($1.95 \times 10^6$ elements) and validated pressure loss and $C_c\%$ at $0$ and $60\text{ RPM}$. \textit{Why cited:} Demonstrates that an Eulerian--Eulerian CFD model with Wen--Yu drag can successfully reproduce the \citet{Osgouei2010} flow-loop trends, directly motivating our geometric setup and boundary conditions. \textit{Critical Limitations:} (a) \citet{Shehab2015} modeled the drilling fluid exclusively as \textbf{Newtonian water} ($\mu = 0.001\text{ Pa}\cdot\text{s}$), completely omitting non-Newtonian Herschel--Bulkley yield stress and shear-thinning behavior; (b) relied on a proprietary commercial solver (ANSYS CFX) with high mesh cell counts due to unstructured meshing; and (c) did not perform progressive decoupled single-phase verification to isolate viscous creeping flow vs.\ yield-stress plug mechanics prior to multiphase coupling.
    
    \item \textbf{\citet{Sorgun2010} (Experimental \& Concentric Modeling Benchmark):} Investigated non-Newtonian water-based muds in eccentric horizontal flow loops and developed a finite-difference MATLAB code to solve turbulent flow with pipe rotation. \textit{Why cited:} Provides realistic Herschel--Bulkley rheological parameters ($K, n, \tau_y$) and experimental insights into how pipe rotation reduces bed thickness in non-Newtonian fluids. \textit{Critical Limitation:} The mathematical finite-difference code developed by \citet{Sorgun2010} was restricted to \textbf{single-phase Newtonian fluids in a concentric annulus}, leaving the 3D multiphase non-Newtonian eccentric problem unresolved in open code.
    
    \item \textbf{\citet{Hajipour2020} \& \citet{Khaled2020} (Commercial Non-Newtonian CFD Benchmarks):} Utilized Eulerian--Eulerian KTGF models in commercial CFD software to evaluate Herschel--Bulkley mud rheology, ROP, eccentricity, and rotation on cuttings concentration. \textit{Why cited:} Validates the theoretical coupling of Herschel--Bulkley rheology with KTGF granular temperature ($\Theta_s$) and Lun et al.\ / Syamlal--O'Brien solid stress closures. \textit{Critical Limitation:} Both studies relied on closed-source commercial solvers and focused primarily on bulk cross-sectional averaged cuttings concentration curves rather than systematically mapping the 2D/3D spatial viscosity field ($\nu_{\text{eff}}$) and classifying the cross-sectional phase regime boundaries (Stationary $\rightarrow$ Moving $\rightarrow$ Dispersed) in an open-source FVM environment.
\end{enumerate}

\begin{table}[htbp]
\centering
\scriptsize
\caption{Critical Synthesis of Key Literature Benchmarks and Comparison with Present Study}
\label{tab:lit_synthesis}
\begin{tabularx}{\textwidth}{l l l l X}
\toprule
\textbf{Author (Year)} & \textbf{Approach / Tool} & \textbf{Fluid Rheology} & \textbf{Key Parameters} & \textbf{Achievements \& Critical Limitations (Research Gap)} \\
\midrule
\citet{Osgouei2010} & Experimental Flow Loop & Water \& Gasified Mud & $D_w=73.9\text{ mm}$, $D_p=47\text{ mm}$, $e=0.623$, $0\text{--}120\text{ RPM}$ & \textbf{Achieved:} Ground-truth $\Delta P/L$, $C_c\%$, and visual flow pattern data. \textbf{Limitation:} No internal 3D velocity, shear rate, or local viscosity field resolution. \\
\addlinespace
\citet{Sorgun2010} & Experimental \& MATLAB FD & Herschel--Bulkley \& Power Law & $e=0.623$, $v_l=0.6\text{--}1.8\text{ m/s}$, $0\text{--}120\text{ RPM}$ & \textbf{Achieved:} Quantified non-Newtonian mud cleaning in flow loop. \textbf{Limitation:} Numerical MATLAB model restricted to concentric Newtonian single-phase flow. \\
\addlinespace
\citet{Shehab2015} & 3D CFD (ANSYS CFX 15.0) & Newtonian Water Only & $e=0.623$, $d_p=2.01\text{ mm}$, $\text{ROP}=60\text{ ft/hr}$, $0\text{--}120\text{ RPM}$ & \textbf{Achieved:} Validated Eulerian CFD against Osgouei (2010) at $0, 60\text{ RPM}$. \textbf{Limitation:} Used water instead of non-Newtonian mud; commercial black-box solver; unstructured $1.95\text{M}$ mesh. \\
\addlinespace
\citet{Hajipour2020} & 3D CFD (Commercial TFM) & Herschel--Bulkley (LVM / HVM) & $e=0\text{--}0.8$, $0\text{--}150\text{ RPM}$, KTGF closures & \textbf{Achieved:} Demonstrated KTGF + Herschel--Bulkley accuracy vs.\ Tomren (1986). \textbf{Limitation:} Commercial closed solver; did not isolate unyielded plug vs.\ narrow-gap shear liquefaction fields. \\
\addlinespace
\citet{Khaled2020} & 3D CFD (Commercial TFM) & Non-Newtonian Muds & Eccentric horizontal annulus, $Y_p/P_v$ ratios & \textbf{Achieved:} Evaluated rheological sensitivity on hole cleaning. \textbf{Limitation:} Lacked open-source verification pipeline and regime transition mapping on Osgouei (2010) geometry. \\
\addlinespace
\citet{Li2026} & Bibliometric \& Critical Review & Multiphase / Non-Newtonian & Cross-scale bed-state variables (625 studies) & \textbf{Achieved:} Identified bridge variables ($\alpha_s$, $\Delta P/L$, ECD) and model fidelity trade-offs. \textbf{Limitation:} Review paper calling for transparent, geometry-matched CFD validation workflows. \\
\addlinespace
\textbf{Present Study (2026)} & \textbf{OpenFOAM v2406 (\texttt{twoPhaseEulerFoam})} & \textbf{Newtonian \& Herschel--Bulkley} & \textbf{$e=0.623$, $d_p=2.01\text{ mm}$, 16-Run Matrix ($v_l \times \Omega$)} & \textbf{Novelty:} Open-source FVM implementation; 3-stage single-phase physical verification (\texttt{stokesCavity}, \texttt{eccentricNS}, \texttt{eccentric\_HB}) + full non-Newtonian KTGF multiphase regime mapping. \\
\bottomrule
\end{tabularx}
\end{table}

\section{Research Gaps and Novelty of the Present Study}
Based on the critical literature synthesis summarized in Table \ref{tab:lit_synthesis}, three distinct research gaps are identified in existing undergraduate and postgraduate cuttings transport studies:
\begin{enumerate}
    \item \textbf{Over-Reliance on Newtonian Carrier Fluids in Validated Eccentric Benchmarks:} While \citet{Shehab2015} successfully modeled the \citet{Osgouei2010} eccentric horizontal wellbore ($e = 0.623$), the simulation was restricted to Newtonian water. Consequently, the interaction between drillpipe rotation ($0\text{--}120\text{ RPM}$) and non-Newtonian Herschel--Bulkley yield-stress rheology ($\tau_y, K, n$)—specifically how rotational shear liquefies the mud near the pipe wall while an unyielded plug persists in the wide gap—was left unexamined in that benchmark geometry.
    \item \textbf{Dependence on Proprietary Commercial ``Black-Box'' CFD Software:} Almost all existing 3D Eulerian--Eulerian wellbore studies \citep{Shehab2015, Hajipour2020, Khaled2020, Ferroudji2022} rely on commercial licenses (ANSYS CFX or Fluent). There is a scarcity of fully documented, transparent, open-source \textbf{OpenFOAM} frameworks that detail the exact kinematic-to-dynamic governing equation transformations, dictionary configurations, and structured O-grid meshing for eccentric horizontal wellbores.
    \item \textbf{Lack of Progressive Single-Phase Verification Prior to Multiphase Execution:} Previous studies jump directly into complex 3D transient multiphase runs on millions of unstructured cells, obscuring numerical discretization errors, CFL instabilities, and non-orthogonal mesh pressure oscillations.
\end{enumerate}

\textbf{Novelty of the Present Work:} To bridge these gaps, this project introduces a rigorous, two-stage open-source CFD methodology in \textbf{OpenFOAM v2406}. First, in FYP 1, we establish a progressive three-method single-phase verification pipeline (\texttt{stokesCavity} $\rightarrow$ \texttt{eccentricNS\_clean} $\rightarrow$ \texttt{eccentric\_HB}) that mathematically and numerically isolates creeping low-side flow ($\text{Co} = 0.20$), eccentric rotational Couette pressure wedges ($60\text{ RPM}$), and non-Newtonian Herschel--Bulkley shear liquefaction vs.\ unyielded plug formation ($\nu_{\text{eff}} \in [1.1 \times 10^{-4}, 1.0 \times 10^{-2}]\text{ m}^2/\text{s}$). Second, in FYP 2, this verified framework is extended into a 3D Eulerian--Eulerian KTGF model (\texttt{twoPhaseEulerFoam}) validated first with Newtonian water against \citet{Osgouei2010} and \citet{Shehab2015}, and then deployed across a 16-run non-Newtonian Herschel--Bulkley parametric matrix ($v_l \in [1.2, 2.4]\text{ m/s}$, $\Omega \in [0, 120]\text{ RPM}$) to map cross-sectional flow regime transitions and optimize hole cleaning efficiency.